%
% 43-fundamentalsatz.tex -- Der Fundamentalsatz der Algebra
%
% (c) 2026 Prof Dr Andreas Müller
%
\subsection{Der Fundamentalsatz der Algebra
\label{buch:analytisch:ganz:subsection:fundamentalsatz}}
Die komplexen Zahlen wurden als Körpererweiterung der reellen Zahlen
konstruiert, um beliebige Quadratwurzeln ziehen zu können.
Damit ist noch nicht klar, dass sich jetzt auch Wurzeln höherer
Ordnung ziehen lassen oder noch allgemeiner, dass jedes Polynom
eine Nullstelle in $\mathbb{C}$ hat.

\begin{satz}[Fundamentalsatz der Algebra]
\label{buch:analytisch:ganz:satz:fundamentalsatz}
\index{Fundamentalsatz der Algebra}%
Das Polynom $f(z) = a_0 + a_1z + \ldots + a_{n-1}z^{n-1} + a_nz^n$ mit
$n>0$ hat mindestens eine komplexe Nullstelle.
\end{satz}

\begin{proof}
Wir untersuchen das Verhalten des Betrags von $f(z)$ für $z\to\infty$.
Dazu schreiben wir
\begin{equation}
|f(z)|
= 
|a_nz^n|
\cdot
\biggl|
\frac{a_0}{a_nz^n}
+
\frac{a_1}{a_nz^{n-1}}
+
\ldots
+
\frac{a_k}{a_nz^{n-k}}
+
\ldots
+
\frac{a_{n-1}}{a_nz}
+
1
\biggr|.
\label{buch:analytisch:ganz:fundamental:eqn:1}
\end{equation}
Wenn $|z|$ gross genug ist, dann ist die $1$ auf der rechten Seite
der dominierende Term.
Wählt man $r$ so gross, dass für alle $k$, $0\le k\le n-1$
\begin{equation}
\frac{a_k}{a_n}
\le
\frac{r^{n-k}}{n+1}
\label{buch:analytisch:ganz:fundamental:proof:3}
\end{equation}
gilt, dann ist jeder Terme auf der rechten Seite von
\eqref{buch:analytisch:ganz:fundamental:eqn:1}
\[
\biggl|
\frac{a_k}{a_nz^{n-k}}
\biggr|
=
\biggl|
\frac{a_k}{a_n}
\biggr|
\cdot
\frac{1}{r^{n-k}}
\stackrel{\eqref{buch:analytisch:ganz:fundamental:proof:3}}{\le}
\frac{r^{n-k}}{n+1}
\cdot
\frac{1}{r^{n-k}}
=
\frac{1}{n+1}.
\]
Aufgelöst nach $r$ bedeutet dies, dass 
\begin{equation}
r\ge
\max
\Bigl\{
\!\sqrt[n-k]{(n+1)|a_k/a_n|}
\;
\Big|
\;
0\le k\le n-1
\Bigr\}
\label{buch:analytisch:ganz:fundamental:proof:r}
\end{equation}
gewählt werden muss.

Mit der Wahl \eqref{buch:analytisch:ganz:fundamental:proof:r}
von $r$ wird die Gleichung
\eqref{buch:analytisch:ganz:fundamental:eqn:1}
für $|z|\ge r$
wegen
\eqref{buch:analytisch:ganz:fundamental:proof:3}
zur Ungleichung
\begin{equation}
|f(z)|
\ge
|a_nr^n|
\cdot
\biggl|
1
-
n
\cdot
\frac{1}{n+1}
\biggr|
=
|a_nr^n|
\cdot
\frac{1}{n+1}.
\label{buch:analytisch:ganz:fundamental:eqn:2}
\end{equation}
Der Betrag von $f(z)$ wächst also mit $r^n$.
Wir haben also gezeigt, dass der Term höchsten Grades im Polynom
das Wachstum dominiert.

Wir nehmen jetzt an, dass $f(z)$ keine Nullstelle hat, und führen
diese Annahme zu einem Widerspruch.
Dann ist $1/f(z)$ eine ganze Funktion.
Aus der Ungleichung
\eqref{buch:analytisch:ganz:fundamental:eqn:2}
folgt dann, dass
\begin{equation}
\frac{1}{|f(z)|}
\le
\frac{n+1}{ |a_n| r^n }
\label{buch:analytisch:ganz:fundamental:eqn:3}
\end{equation}
für alle $|z|\ge r$ sein muss.
$1/f(z)$ ist innerhalb des Kreises vom Radius $r$ sicher beschränkt
und wegen
\eqref{buch:analytisch:ganz:fundamental:eqn:3}
auch ausserhalb.
Nach dem Satz von Liouville ist $1/f(z)$ daher eine Konstante und
daher auch $f(z)$.
Insbesondere kann $f(z)$ kein Polynom vom Grad $\ge 1$ sein.
Dieser Widerspruch zeigt, dass $f(z)$ eine Nullstelle haben muss.
\end{proof}

Der Satz~\ref{buch:analytisch:ganz:satz:fundamentalsatz}
verspricht nur die Existenz einer Nullstelle.
Mithilfe des Divisionsalgorithmus für Polynome lässt sich
daraus sofort eine Faktorisierung in Linearfaktoren konstruieren.
Dies wird im folgenden Satz gezeigt.

\begin{satz}
Jedes Polynom $f(z) = a_0 + a_1z + \ldots + a_{n-1}z^{n-1} + a_nz^n$
vom Grad $n$ mit Koeffizienten in $\mathbb{C}$ lässt sich in $n$
Linearfaktoren
\[
f(z)
=
a_n
(z - z_1)
\cdot
\ldots
\cdot
(z - z_n)
\]
zerlegen, wobei die Zahlen $z_k\in\mathbb{C}$ die Nullstellen von
$f(z)$ sind.
\end{satz}

\begin{proof}
Wir beweisen die Aussage mithilfe vollständiger Induktion.
Ein Polynom ersten Grades ist bereits ein Linearfaktor, dies ist die
Induktionsverankerung.

Zur Verankerung der Induktion nehmen wir jetzt an, dass bereits gezeigt
ist, dass alle Polynome vom Grad $n-1$ sich als Produkt von $n-1$
Linearfaktoren schreiben lassen.

Das Polynom $f(z)$ hat nach
Satz~\ref{buch:analytisch:ganz:satz:fundamentalsatz} eine
Nullstelle $z_n$.
Der Divisionsalgorithmus für Polynome liefert dann Polynome
$q(z)$ und $r(z)$ derart, dass
\[
f(z) = (z-z_n)\cdot q(z) + r(z).
\]
Der Grad von $r(z)$ muss kleiner sein als der Grad des Divisors $z-z_n$.
Somit hat $r(z)$ Grad $0$ und ist eine Konstante.
Da $f(z_n)=0$ und $z-z_n=0$ ist, muss auch $r(z_n)=0$ sein, daher
ist die Konstante $0$.

Das Polynom $q(z)$ hat Grad $n-1$ und ist daher nach Induktionsannahme
ein Produkt
\[
q(z)
=
a_n(z-z_1)\cdot\ldots\cdot (z-z_{n-1})
\]
von Linearfaktoren.
Also ist auch $f(z)$ ein Produkt von Linearfaktoren.
\end{proof}

Der Fundamentalsatz sagt nicht, dass es für die Nullstellen eines
Polynoms eine einfach Formel gibt, wie dies zum Beispiel für die
Nullstellen
\[
x_{1,2}
=
\frac{-a_1\pm\!\sqrt{a_1^2-4a_0a_2}}{2a_2}
\]
eines quadratischen Polynoms der Fall ist.
Auch für kubische Polynome gibt es eine solche Formel.
Das normierte kubische Polynom $x^3+a_2x^2+a_1x+a_0$ 
kann durch die Verschiebung $u = x-a_2/3$ 
in die Form
\[
u^3
+
\biggl(
\underbrace{
\frac13 a_2^2
-\frac23 a_1a_2
+a_1
}_{\displaystyle = p}
\biggr)
u
+
\biggl(
\underbrace{
\frac{1}{27} a_2
+
\frac19 a_2^3
-\frac13 a_1a_2
+
a_0
}_{\displaystyle = q}
\biggr)
=
u^3 + pu + q
\]
gebracht werden.
Eine Nullstelle dieses Polynoms ist dann durch die cardanische Formel
\index{cardanische Formel}%
\[
u
=
\!\sqrt[\scriptstyle3]{
-\frac{q}3+\!\sqrt{
\rlap{$\raisebox{0.5pt}{\phantom{\bigg|}}$}
\smash{
\biggl(\frac{q}2\biggr)^{\raisebox{-2pt}{$\scriptstyle 2$}}
+
\biggl(\frac{p}3\biggr)^{\raisebox{-2pt}{$\scriptstyle 3$}}
}}
}
+
\!\sqrt[\scriptstyle3]{
-\frac{q}3-\!\sqrt{
\rlap{$\raisebox{0.5pt}{\phantom{\bigg|}}$}
\smash{
\biggl(\frac{q}2\biggr)^{\raisebox{-2pt}{$\scriptstyle 2$}}
+
\biggl(\frac{p}3\biggr)^{\raisebox{-2pt}{$\scriptstyle 3$}}
}}
}
\]
gegeben.

Auch für Polynome vierten Grades gibt es ähnliche Formeln 
\index{Polynom vierten Grades}%
von Lodovico Ferrari.
\index{Ferrari, Lodovico}%
Niels Henrik Abel hat jedoch gezeigt, dass es für $n>4$ im
\index{Abel, Niels Henrik}%
Allgemeinen keine solche Formel gibt.
Es ist zwar wahr, dass es für jedes Polynom vom Grad $n$
über $\mathbb{C}$ eine Produktdarstellung aus Linearfaktoren
gibt, aber es gibt kein algebraische Methode, diese Nullstellen in
Form eines algebraischen Ausdrucks zu finden.
Mit dem Computer sind beliebig genaue numerische Lösungen
zugänglich.
Für ein Computeralgebrasystem entsteht jedoch die Schwierigkeit,
dass es keinen Algorithmus gibt, diese Nullstellen exakt zu finden.

Für die Integration rationaler Funktionen wird zum Beispiel gelehrt,
dass die Partialbruchzerlegung gebildet werden muss, für deren einzelne
\index{Partialbruchzerlegung}%
Summanden Stammfunktionen relativ einfach gefunden werden können.
Die Bildung der Partialbruchzerlegung verlangt aber, dass die
Nullstellen des Nenners in algebraischer Form gefunden werden
müssen.
Da es dafür keinen Algorithmus gibt, kann ein Computeralgebrasystem
diesen Algorithmus nur in Spezialfällen implementieren, in denen
eine Faktorisierung des Nenners bekannt ist.
Diese Schwierigkeit wird in den Kapiteln~\ref{chapter:ratint}
und \ref{chapter:intalg} adressiert.


